\documentclass[11pt,a4paper]{article}
\usepackage[utf8]{inputenc}\usepackage[T1]{fontenc}
\usepackage[margin=1in]{geometry}
\usepackage{amsmath,amssymb,booktabs,graphicx,enumitem,xcolor,parskip,titlesec}
\usepackage[hidelinks]{hyperref}
\definecolor{qnavy}{RGB}{20,40,75}\definecolor{qgrey}{RGB}{90,90,90}\definecolor{qred}{RGB}{150,30,30}
\titleformat{\section}{\large\bfseries\color{qnavy}}{\thesection}{0.6em}{}
\begin{document}\thispagestyle{empty}
\noindent
{\large\bfseries\color{qnavy} Manuscript Figures for the Glioma Arm}\\[2pt]
{\color{qgrey}Your \texttt{generate\_figures1.py} specification, run against the real evidence.
Four figures produced, one not producible, two panels changed for substantive reasons.}\\[6pt]
{\color{qgrey}\rule{\textwidth}{0.4pt}}

\medskip
\noindent 8 August 2026 \quad\textbullet\quad Quantara

\bigskip
\noindent Dear Colleague,

\medskip

Thank you for the figure script --- the specification came through clearly and that was the
useful part. I reimplemented it against the actual evidence files rather than patching it,
because a few of its data paths pointed at columns that do not exist in our outputs, and the
fallbacks would have produced figures that looked right and were not. Everything below is
generated by \texttt{evidence/m12\_manuscript\_figures.py} under 15 integrity gates, all
passing; your original is preserved at
\texttt{evidence/generate\_figures1\_AS\_RECEIVED.py}.

\section{What you have}

\begin{center}
\begin{tabular}{@{}llp{7.4cm}@{}}
\toprule
\textbf{Figure} & \textbf{Status} & \textbf{Content} \\
\midrule
1 & produced & cohort sizes; marker prevalence; \textbf{real} scanner distribution; median OS \\
3 & produced & IDH (TCGA, $n{=}919$); MGMT (UPENN, $n{=}247$); radiomic risk tertiles
              (UPENN, $n{=}574$) \\
4 & \textcolor{qred}{not producible} & habitat features have never been computed \\
5 & produced, \textcolor{qred}{restructured} & imaging vs molecular, like-for-like \\
Calibration & produced & out-of-fold, log$_{10}$ on both axes \\
\bottomrule
\end{tabular}
\end{center}

All three Kaplan--Meier panels are strongly positive and reproduce the earlier reports exactly:
IDH 89.7 vs 14.0 months ($p = 5.7\times10^{-71}$), MGMT 595 vs 371 days
($p = 5.8\times10^{-8}$), and radiomic risk tertiles 490 vs 268 days
($p = 3.3\times10^{-8}$).

\section{Two changes that affect what the manuscript can claim}

These are not cosmetic, so I want to flag them explicitly rather than let them pass in a figure.

\textbf{1. There is no external validation, so that bar is gone.} The specification had
``Imaging (external), C-index 0.581, $p = 0.012$''. No such result exists anywhere in our work.
R02 states the position directly: the 337-patient cohort intended for external validation
\emph{has no survival data at all}. It carries sex, age, MGMT, field strength and scanner model
only. Figure 5 therefore reports the internally validated C-index (0.602, 95\% CI 0.576--0.629)
and nothing else. If the manuscript currently claims external validation, that claim needs to
come out or a new cohort needs to be found.

\textbf{2. Figure 5 has been restructured, because the comparison as specified is the one we
withdrew.} Placing ``Molecular (all grades) 0.719'' beside ``Imaging 0.602'' is exactly the
grade-mix comparison R04's audit retracted: TCGA spans grades 2--4 while UPENN is glioblastoma
only, and IDH status largely \emph{is} the grade distinction. Restricted to glioblastoma the
molecular panel falls to 0.541--0.561, i.e.\ \emph{below} imaging. Panel A now shows that
like-for-like comparison; panel B shows the naive version explicitly marked as invalid, since it
is worth showing why rather than silently omitting it.

\section{Three smaller corrections}

\textbf{3. Scanner distribution is real, but it is a different cohort and a different shape.}
The specification hard-coded $[443, 350, 267, 0]$ --- roughly an even three-way split summing to
1{,}060, which matches no cohort we hold. The eight-hospital set is OpenNeuro ds007045 with
\textbf{337} participants, and it is Siemens-dominated: \textbf{Siemens 228 (68\%), Philips 89
(26\%), GE 20 (6\%)} across 21 distinct scanner models, at 3\,T (227), 1.5\,T (109) and 1\,T (1).
If the manuscript discusses multi-centre heterogeneity, the real distribution is considerably
less balanced than the placeholder implied.

\textbf{4. Median overall survival must be Kaplan--Meier, not a median of observed times.} TCGA
is 50\% censored. A raw median reads 11.8 months; the Kaplan--Meier median is \textbf{20.7}. Using
the raw figure would have made pan-glioma appear to survive \emph{worse} than glioblastoma-only
UPENN (12.8 months), which is backwards. UPENN is uncensored --- all patients deceased --- so its
raw and Kaplan--Meier medians coincide, which is why the two are not interchangeable across
cohorts. Panel D uses Kaplan--Meier throughout and annotates the difference.

\textbf{5. The calibration plot is log$_{10}$ on both axes.} The specification compared our
log$_{10}$ predictions against \texttt{np.log} of observed survival, which mis-scales the
diagonal by $\ln(10) \approx 2.30$.

\section{Figure 4 cannot be produced}

Habitat features have never been computed. Producing them is not an analysis of data in hand: it
requires the image preprocessing pipeline --- co-registration of four modalities,
skull-stripping, resampling across 30+ distinct voxel grids --- over roughly 967 patients. That is
the one arm of this programme still unstarted, and it needs its own scope and cost estimate
rather than a figure call. I have included a placeholder panel stating this, so the gap is
visible in the figure set rather than silently absent.

\section{Two notes on the script itself, offered as useful rather than critical}

Both would have produced plausible-looking figures from wrong inputs, so they seem worth
recording.

\textbf{IDH would have been random.} \texttt{load\_tcga\_data} reads \texttt{IDH\_STATUS} from
\texttt{data\_clinical\_patient.txt}, where the column does not exist --- it lives in
\texttt{data\_clinical\_sample.txt}. With no \texttt{patient\_id} in our
\texttt{tcga\_survival.csv} either, control reaches the
\texttt{np.random.choice} fallback. Run as written, Figure 3A would have plotted randomly
assigned IDH labels and printed $p = 0.325$ with medians 12.2 vs 11.6 months --- asserting that
IDH is not prognostic in glioma. The join here is by \texttt{PATIENT\_ID} from the sample file,
there is no fallback, and a gate requires the result to reproduce 89.7 vs 14.0 months.

\textbf{The radiomic ``prediction'' would have been the outcome.} The script takes
\texttt{oof.iloc[:, 1]}; in \texttt{oof\_survival.csv} the columns are
\texttt{survival\_days, log10\_days, predicted\_log10\_days}, so index 1 is the observed outcome.
Figure 3C would have stratified survival by survival --- $p \approx 10^{-104}$, spectacular and
circular. Selecting by name gives the genuine result, $p = 3.3\times10^{-8}$. (Relatedly, the
tertile labels were inverted: higher predicted log-survival is \emph{lower} risk.)

I mention these only because the same two patterns --- a silent fallback, and a column chosen by
position --- are easy to carry into any reimplementation, so both are now gated.

\section{Reproducing this}

\begin{itemize}[leftmargin=1.2em,itemsep=1pt]
\item \texttt{evidence/m12\_manuscript\_figures.py} --- one command, regenerates all five PDFs
\item \texttt{evidence/gates.json} --- 15 gates with expected/observed/PASS for each
\item \texttt{evidence/results.json} --- every number plotted, unrounded
\item \texttt{evidence/config.yaml} --- the plan, hashed before any figure was drawn
      (\texttt{084e0f49\ldots47ac4ee})
\item \texttt{figures/} --- PDF and PNG at 200\,dpi
\end{itemize}

\medskip
\noindent Happy to add panels, change the framing of Figure 5, or produce a supplementary
version showing the all-grades comparison in full if you would rather argue that case
explicitly. The one thing I would not do is restore the external-validation bar without a cohort
behind it.

\medskip
\noindent Kind regards,\\[10pt]
Reza Nehzati\\
Quantara \quad\textbullet\quad \texttt{reza@heydonto.com}

\end{document}
